% D.71. Root and PBC independent full proof and primary-source audits PASS.
\subsection{Generic class-two rigidity on the once-punctured elliptic curve}

\begin{statement}{D.71}{Proved}{The requested one-cusp curve has generic singleton class-two fibres}
Let \(E/k\) be an elliptic curve over a number field, let
\(U=E-\{O\}\), and fix a nonzero rational tangent at \(O\).
Use the resulting actual cusp section \(s_0\) as base point.
For every prime \(\ell\), put
\[
 P_\ell=\pi_1(U_{\bar k})^{(\ell)},\qquad
 T_{2,\ell}=P_\ell/\gamma_3P_\ell,\qquad
 V_\ell=T_\ell E,\qquad
 {\cal K}_{k,\ell}=H^1(G_k,\mathbb Z_\ell(1)).
                                                        \tag{D.71.1}
\]
The actual cusp packet gives a bijection
\[
 {\cal K}_{k,\ell}\ \xrightarrow{\ \sim\ }\
 \ker\!\left(H^1(G_k,T_{2,\ell})\longrightarrow
                                  H^1(G_k,V_\ell)\right).
                                                        \tag{D.71.2}
\]
For a dense \(G_\delta\) set of parameters in the intrinsic
\(\ell\)-adic topology on \({\cal K}_{k,\ell}\), the full
pro-\(\ell\) section fibre above its class in (D.71.2)
is a singleton, represented by that actual completed cusp.
The parameters may all be chosen infinitely ramified.

The same statement holds for the actual full profinite
class-two quotient and a dense \(G_\delta\) set in
\(\widehat{\cal K}_k=\prod_\ell{\cal K}_{k,\ell}\):
its full profinite section fibre is the singleton cusp
class. All primewise parameters may simultaneously
be infinitely ramified. This conclusion uses precisely
D.24 and D.56.

In particular these assertions hold for the specified
\[
 k=\mathbb Q(i),\qquad E:y^2=x^3-x,\qquad U=E-\{O\}.
                                                        \tag{D.71.3}
\]
Every cusp parameter has the same abelian section
class, but distinct parameters already have distinct
class-two classes. Thus the abelian fibre itself
contains an injected copy of the entire completed
cusp packet and cannot be a singleton. The new
generic uniqueness is at class two. No first finite
nilpotent or derived obstruction is asserted.
\end{statement}

\begin{proof}
\paragraph{The actual integral class-two group.}
Fix \(\ell\), write \(P=P_\ell\), and let \(c\)
generate the actual cusp inertia. The group \(P\)
is free pro-\(\ell\) on two generators. In a
surface basis the peripheral word is a commutator
or its inverse, according to orientation.
Consequently its class in
\(\gamma_2P/\gamma_3P\) is a primitive generator.
The commutator pairing and the Weil pairing give
the integral equivariant identifications
\[
 P^{\rm ab}=V=T_\ell E,\qquad
 I:=\gamma_2P/\gamma_3P
       =\bigwedge^2 V\simeq\mathbb Z_\ell(1).
                                                        \tag{D.71.4}
\]
In particular these statements hold at two; no
Baker--Campbell--Hausdorff denominator is used.
The actual central extension is
\[
 1\longrightarrow I\longrightarrow T_2
                     \longrightarrow V\longrightarrow1.
                                                        \tag{D.71.5}
\]
The map \(\langle c\rangle\to I\) is an isomorphism
of \(\mathbb Z_\ell(1)\)-modules. The tangential
action \(\sigma_0\) satisfies
\(\sigma_0(g)c=c^{\chi_\ell(g)}\).

Choose a finite set \(S\) containing all infinite
places, all places above \(\ell\), the bad
reduction places of \(E\), and the places where
the fixed tangent is not an integral nonzero
tangent of the smooth model. At \(v\notin S\),
\[
             V^{G_{k_v}}=0,\qquad I^{G_{k_v}}=0.
                                                        \tag{D.71.6}
\]
The first assertion follows from injection of
local prime-to-residue-characteristic torsion
into the finite reduction group, exactly as
proved in D.70. The second follows because
Frobenius acts on \(I\) by \(N(v)\ne1\).
Restriction to one such place gives the same
global vanishings. A fixed element of \(T_2\)
has zero image in \(V\) and then lies in
\(I^{G_{k_v}}\), or in \(I^{G_k}\), respectively.
Thus
\[
             T_2^{G_{k_v}}=T_2^{G_k}=1
                         \qquad(v\notin S).             \tag{D.71.7}
\]
These are vanishings for the actual integral
groups.

Any \(T_2\)-cocycle with zero abelian class
can be conjugated to have zero actual abelian
cocycle, by lifting a conjugator from \(V\).
It is then an \(I\)-valued cocycle. A conjugator
between two such representatives has
Galois-fixed image in \(V\), which is zero.
It therefore belongs to \(I\), and its
equivalence relation is exactly the usual
integral \(I\)-coboundary relation. This
proves (D.71.2), globally and also locally
at every \(v\notin S\). Every parameter is
supplied by the actual cusp cocycle
\(g\mapsto c^{a(g)}\).

\paragraph{The centralizer of the nonprimitive peripheral word.}
We prove the required free pro-\(\ell\) fact,
including its subgroup input. A finitely
generated closed subgroup \(L\) of a finite
rank free pro-\(\ell\) group \(P\) is a
free factor of some open subgroup. This
is \cite[Theorem 3.2, pp.~317--318]{DE_Lubotzky1982};
here is its finite-character proof in the
form needed here.

Choose a basis of the finite vector space
\(\operatorname{Hom}_{\rm cts}(L,\mathbb F_\ell)\).
For a common sufficiently small open normal
subgroup \(W\lhd P\), every chosen character
vanishes on \(L\cap W\). Each extends to
the open subgroup \(LW\) by
\(\phi(lw)=\phi(l)\). This is well-defined
and a homomorphism because \(W\) is normal.
Thus \(L/\Phi(L)\to LW/\Phi(LW)\) is injective.
The group \(LW\) is finite rank free pro-\(\ell\),
by the open-subgroup Schreier theorem.
A minimal generating set of \(L\) therefore
extends to a minimal generating set of \(LW\).
Such a set is a free basis: the resulting
surjective endomorphism of a finite rank free
pro-\(\ell\) group is an isomorphism. The
Hopf property follows directly because a
finitely generated profinite group has only
finitely many open subgroups of each index;
a surjective endomorphism permutes their
finite intersections and is injective on
the inverse limit of the resulting finite
quotients. The generators of \(L\) are
therefore part of a free basis, as claimed.

The peripheral element \(c\) is not an
\(\ell\)-th power in \(P\). If \(c=t^\ell\),
its zero abelian image forces \(t\in\gamma_2P\),
since \(V\) is torsion-free. Its primitive
class in \(\gamma_2P/\gamma_3P\) would then
be divisible by \(\ell\), a contradiction.
For \(0\ne b\in\mathbb Z_\ell\) and
\(g\in C_P(c^b)\), consider the closed
two-generated subgroup \(L=\langle g,c\rangle\).
The preceding lemma makes it free pro-\(\ell\).
If it had rank two, \(g,c\) would be a
free basis and \([g,c^b]\ne1\). One sees
the latter directly in the Magnus embedding:
the coefficient \(b\) of the degree-two
word \(GC-CG\) is nonzero in
\(\mathbb Z_\ell\). Hence \(L\) is procyclic.
As \(c\) has no \(\ell\)-th root in \(P\),
it generates \(L\). We obtain
\[
                   C_P(c^b)=\langle c\rangle
                         \quad(b\ne0).                 \tag{D.71.8}
\]
This proof does not treat \(c\) as a primitive
free generator.

The geometric normalizer of any nonzero
closed subgroup of \(\langle c\rangle\)
also equals \(\langle c\rangle\).
Indeed, if conjugation sends \(c^b\)
to \(c^{bu}\), its effect on
\(\gamma_2P/\gamma_3P\) is trivial.
Thus \(b[c]=bu[c]\); the primitive
torsion-free class forces \(u=1\),
and (D.71.8) applies.
Finally, for \(H=\gamma_3P\),
\[
                 H\cap\langle c\rangle=1.              \tag{D.71.9}
\]

\paragraph{The exact fibre and its local comparison.}
Use the actual Kummer representatives
on each fixed torsion component of
\({\cal K}_{k,\ell}\) constructed in
D.70. Define
\[
       \sigma_a(g)=\operatorname{Int}(c^{a(g)})\sigma_0(g)
                         \quad\hbox{on }P\hbox{ and }H.
                                                        \tag{D.71.10}
\]
On \(T_2\) this action equals \(\sigma_0\),
because the image of \(c\) is central.
A full pro-\(\ell\) section with the
same class-two class as this cusp can
be conjugated to have exactly its
class-two cocycle. Its difference is
then a continuous \(H\)-cocycle.
Conversely every such cocycle gives
a section in the fibre. A conjugator
between exact representatives projects
to a fixed element of \(T_2\); (D.71.7)
makes it trivial. Thus the exact fibre is
\[
                         H^1(G_k,H;\sigma_a).          \tag{D.71.11}
\]

At \(v\notin S\), apply the cartesian
good-reduction specialization square
and the inertia alternative of
\cite[Definition A.3, Remark A.3.1 and Proposition A.6]
{DE_Hoshi2012} with \(\Sigma=\{\ell\}\).
There is only one geometric cusp.
If \(a_v\ne0\), the prescribed
class-two inertia \(c^{a_v t}\)
is nonzero. The actual section
therefore has nonzero specialized
inertia in that cusp. Its whole
image normalizes this inertia;
(D.71.8) and its normalizer
consequence, compared with the
tangential arithmetic lift, place
the image in the cusp decomposition
group. The cartesian square puts
the actual local section in the
actual local pro-\(\ell\) packet.
The local version of (D.71.2)
shows that its Kummer parameter
is exactly \(a\) as a cohomology
class. Adjusting the packet
conjugator by a cusp-inertia
element makes the two class-two
cocycles agree exactly. Its
projection to \(T_2\) is then
fixed, hence zero by (D.71.7);
the conjugator lies in \(H\).

If \(a_v=0\), nonzero specialized
inertia in the sole cusp would
have a nonzero image in \(I\),
contrary to the prescribed zero
class-two inertia. Both actual
sections therefore have trivial
specialized inertia. For every
cocycle \(z:G_k\to H\) we have
the exact assertions
\[
 \begin{cases}
 z|_{G_{k_v}}\text{ is an }H\text{-coboundary},
                                  &a_v\ne0,\\
 z(I_v)=1,                        &a_v=0,
 \end{cases}\qquad v\notin S.                         \tag{D.71.12}
\]

\paragraph{The Baire construction in the genuine fibre.}
All parameter-space and arithmetic
constructions in D.70 now apply,
with the following precise
substitutions. Use the same complete
space
\(\mu(k)(\ell)\times c_0({\cal I},\mathbb Z_\ell)\),
with fixed actual torsion cocycles.
Use \(P_m=\Phi^m(P)\), \(Q_m=P/P_m\),
and now
\[
 H_m=\operatorname{im}(\gamma_3P\to Q_m),\qquad
 S_m(a)=S\cup\{v:a_v\notin\ell^m\mathbb Z_\ell\}.
                                                        \tag{D.71.13}
\]
The exponent of \(Q_m\) divides
\(\ell^m\). By (D.71.12), every
integral \(H\)-cocycle reduces to
the finite set
\[
                 Z^1(G_{k,S_m(a)},H_m;\rho_{a,m}).
                                                        \tag{D.71.14}
\]
In particular this finite capture
includes all unrestricted global
cocycles. At a later nonzero
inertia place, its integral local
conjugation coboundary is trivial
on \(Q_m\) because the exponent
of \(c\) is divisible by
\(\ell^m\); at a zero inertia
place use the second line of
(D.71.12). The finite action
and support are constant on
each fixed torsion component
and free-parameter coset
modulo \(\ell^m\).

Impose, for each finite cocycle
nontrivial on the action kernel,
a witness place with identity
Frobenius action, nontrivial
cocycle value, and nonzero
\(a_v\). These conditions are
open. They are dense by the
actual finite affine-image
Chebotarev argument of D.70:
an element in the action kernel
with nontrivial cocycle value
first gives an element of the
actual finite affine image.
Choose its Frobenius place
outside all current finite
supports, and perturb by
\(\ell^M\kappa(\alpha_v)\),
where \(\operatorname{div}(\alpha_v)=hv\).
Avoid valuations not divisible
by \(\ell^{M+v_\ell(h)+1}\).
The new nonzero valuation is
exactly \(M+v_\ell(h)\), while
the finite action and prescribed
neighborhood remain unchanged.
This works also on every torsion
component at two.

Intersect these countably many
open dense conditions and the
complements of all finite-support
loci. The Baire theorem gives
the asserted dense \(G_\delta\)
of infinitely ramified parameters.
For such a parameter, (D.71.12)
excludes every exceptional finite
cocycle at its witness place.
Every actual \(H\)-cocycle
therefore descends continuously
to the actual compact image
\[
 {\cal I}_a=\operatorname{im}\left(
       G_k\longrightarrow\prod_m\operatorname{Aut}(H_m)\right).
                                                        \tag{D.71.15}
\]
Choose \(v\notin S\) with \(a_v\ne0\).
Its inertia image is the normal
subgroup
\[
 \{\operatorname{Int}(c^{a_v t})|_H:t\in\mathbb Z_\ell\}
                    \ \lhd\ {\cal I}_a.                \tag{D.71.16}
\]
Normality follows from
\(\sigma_a(g)c=c^{\chi_\ell(g)}\).
Its fixed subgroup in \(H\)
is trivial by (D.71.8)--(D.71.9),
even for nonprimitive \(a_v\).
The local coboundary in (D.71.12)
allows one conjugation by an
element of \(H\) to kill the
cocycle on (D.71.16). Comparing
the two cocycle formulas at
\(ng=g(g^{-1}ng)\) then makes
every value fixed by (D.71.16);
every value is one. This proves
the singleton in (D.71.11).

\paragraph{Full profinite class two and the abelian obstruction.}
For the actual full geometric
group \(\pi=\pi_1(U_{\bar k})\),
\[
             \pi/\gamma_3\pi\simeq\prod_\ell T_{2,\ell}.
                                                        \tag{D.71.17}
\]
Indeed every finite class-two
quotient is nilpotent and is
the product of its Sylow groups.
The universal finite quotient
property and compactness give
the displayed isomorphism,
with the actual arithmetic
action on every factor.
Continuous cocycles and their
conjugacy relations into this
product are componentwise.
Thus (D.71.2) also holds with
full coefficients
\(\widehat{\mathbb Z}(1)\)
and parameter space
\(\widehat{\cal K}_k\).

Take the product of the
primewise dense \(G_\delta\)
sets just proved, with the
complete product topology
specified in D.70. The product
is dense \(G_\delta\), and each
parameter gives an actual
full cusp cocycle. Any actual
full section with that class-two
class has the same full
pro-\(\ell\) class as this
cusp at every \(\ell\).
D.24, with D.56 as its
proved local input, identifies
the two full section classes.
This proves the full assertion.

Finally \(c\in\gamma_2P\), and
the analogous full peripheral
element is in \(\gamma_2\pi\).
All cusp parameters therefore
have zero abelian cocycle.
Equation (D.71.2) and its
full version make distinct
parameters distinct already
at class two. The abelian
fibre contains the whole
injected packet, as claimed.
This applies without a change
of field or curve to (D.71.3).
\end{proof}
